% Part: first-order-logic
% Chapter: natural-deduction
% Section: propositional-rules

\documentclass[../../../include/open-logic-section]{subfiles}

\begin{document}

\iftag{FOL}
      {\olfileid{fol}{ntd}{prl}}
      {\olfileid{pl}{ntd}{prl}}

\olsection{વિધાનાત્મક નિયમો}


\subsection{$\land$ માટેના નિયમો}

\begin{defish}
\AxiomC{$!A$}
\AxiomC{$!B$}
\RightLabel{\Intro{\land}}
\BinaryInfC{$!A \land !B$}
\DisplayProof
\hfill
\begin{tabular}{r}
\AxiomC{$!A \land !B$}
\RightLabel{\Elim{\land}}
\UnaryInfC{$!A$}
\DisplayProof
\\[3ex]
\AxiomC{$!A \land !B$}
\RightLabel{\Elim{\land}}
\UnaryInfC{$!B$}
\DisplayProof
\end{tabular}
\end{defish}

\subsection{$\lor$ માટેના નિયમો}

\begin{defish}
\begin{tabular}{r}
\AxiomC{$!A$}
\RightLabel{\Intro{\lor}}
\UnaryInfC{$!A \lor !B$}
\DisplayProof
\\[3ex]
\AxiomC{$!B$}
\RightLabel{\Intro{\lor}}
\UnaryInfC{$!A \lor !B$}
\DisplayProof
\end{tabular}
\hfill
\AxiomC{$!A \lor !B$}
\AxiomC{$\Discharge{!A}{n}$}
\DeduceC{$!C$}
\AxiomC{$\Discharge{!B}{n}$}
\DeduceC{$!C$}
\DischargeRule{\Elim{\lor}}{n}
\TrinaryInfC{$!C$}
\DisplayProof
\end{defish}

\subsection{$\lif$ માટેના નિયમો}

\begin{defish}
\AxiomC{$\Discharge{!A}{n}$}
\DeduceC{$!B$}
\DischargeRule{\Intro{\lif}}{n}
\UnaryInfC{$!A \lif !B$}
\DisplayProof
\hfill
\AxiomC{$!A \lif !B$}
\AxiomC{$!A$}
\RightLabel{\Elim{\lif}}
\BinaryInfC{$!B$}
\DisplayProof
\end{defish}

\subsection{$\lnot$ માટેના નિયમો}

\begin{defish}
\AxiomC{$\Discharge{!A}{n}$}
\noLine
\DeduceC{$\lfalse$}
\DischargeRule{\Intro{\lnot}}{n}
\UnaryInfC{$\lnot !A$}
\DisplayProof
\hfill
\AxiomC{$\lnot !A$}
\AxiomC{$!A$}
\RightLabel{\Elim{\lnot}}
\BinaryInfC{$\lfalse$}
\DisplayProof
\end{defish}

\subsection{$\lfalse$ માટેના નિયમો}

\begin{defish}
\AxiomC{$\lfalse$}
\RightLabel{\FalseInt}
\UnaryInfC{$!A$}
\DisplayProof
\hfill
\AxiomC{$\Discharge{\lnot !A}{n}$}
\DeduceC{$\lfalse$}
\DischargeRule{\FalseCl}{n}
\UnaryInfC{$!A$}
\DisplayProof
\end{defish}

$\Intro{\lnot}$ અને $\FalseCl$ બહુ સરખા છે. તફાવત એટલો છે કે
$\Intro{\lnot}$થી નિષેધિત !!{sentence}~$\lnot !A$ નિષ્પન્ન થાય છે,
જ્યારે $\FalseCl$થી વિધેયાત્મક !!{sentence}~$!A$ નિષ્પન્ન થાય છે.

કોઈ નિયમ અમુક ધારણાને નિવૃત્ત કરી શકાય એમ બતાવે ત્યારે આપણે તેને છૂટ
ગણીએ છીએ, આવશ્યકતા નહીં. દાખલા તરીકે, $\Intro{\lif}$ નિયમમાં
આધારવાક્ય~$!B$ની !!{derivation}માં આવતી સ્વરૂપ~$!A$ની શૂન્ય સહિત
કોઈ પણ સંખ્યાની ધારણાઓ નિવૃત્ત કરી શકાય છે.

\end{document}
